% Slide 6 — Feature Extraction Methods
\begin{frame}[shrink=6]{Feature Extraction Methods}
  \begin{columns}[T,onlytextwidth]
    \column{0.49\textwidth}
    \begin{block}{A) Statistical features (intensity)}
      \centering
      \begin{tikzpicture}[x=1cm,y=1cm]
        \draw[->,gray] (0,0) -- (3.1,0) node[below] {\scriptsize intensity};
        \draw[->,gray] (0,0) -- (0,1.6) node[left] {\scriptsize count};
        \draw[accent,thick] plot[smooth] coordinates {(0.2,0.1) (0.6,0.4) (1.1,1.0) (1.6,1.2) (2.1,0.7) (2.6,0.3) (2.9,0.15)};
      \end{tikzpicture}
      \vspace{-0.2em}

      \scriptsize Compact summary of intensity distribution: mean, std, entropy,
      skew/kurtosis, and percentiles (p10--p90).
    \end{block}

    \begin{block}{B) LBP (local texture)}
      \centering
      \includegraphics[width=\linewidth,height=0.19\textheight,keepaspectratio,trim=345 250 340 360,clip]{descriptor_demo.png}
      \vspace{-0.2em}

      \scriptsize Encodes local micro-patterns; robust texture fingerprint for tumor tissue.
    \end{block}

    \column{0.49\textwidth}
    \begin{block}{C) DWT (frequency sub-bands)}
      \centering
      \includegraphics[width=\linewidth,height=0.19\textheight,keepaspectratio,trim=60 0 60 585,clip]{descriptor_demo.png}
      \vspace{-0.2em}

      \scriptsize Multi-scale decomposition (LL/LH/HL/HH) captures coarse vs. fine structures.
    \end{block}

    \begin{block}{D) HOG (edges \& shape)}
      \centering
      \includegraphics[width=\linewidth,height=0.19\textheight,keepaspectratio,trim=685 250 60 360,clip]{descriptor_demo.png}
      \vspace{-0.2em}

      \scriptsize Aggregates gradient orientations; strong for boundary/shape cues.
    \end{block}
  \end{columns}

  \vspace{-0.2em}
  \centering
  {\scriptsize Four complementary views of the same slice: intensity, texture, frequency, and gradients.}
\end{frame}

\note{
Each family captures a different aspect of tumor appearance.
They are later fused after individual optimization.
}
